\section{A complete coherent local detector and record join}
\label{sec:new-join}
The two-arm evolution just derived is an actual source--continuum
construction for the stated receiving action. It does not assume that
every arm pair produces the quarter-cycle age separation used in the
decoder calibration. We give a finite, explicitly paid route-label
encoding that supplies a complete worked experiment, while retaining
the unforced proper-age comparison as a distinct prediction.

\subsection{Closed physical paths and a controlled source encoding}
Take mirror paths with
\[\begin{gathered}
v_{0,1}(t)=\pm v_*\sin(2\pi t/T)\widehat e,\\
 x_{0,1}(t)=\pm\frac{v_*T}{2\pi}[1-\cos(2\pi t/T)]\widehat e,
 \\
 v_*=1/64,\\
 T=2^{-20}.
\end{gathered}\]
Both start and return to the same centre at rest. Their proper ages are
equal. They still couple at different positions to one common field.
Their total variation gives
\(B(T)\leq8v_*+4v_*^2=129/1024<1\), hence
\(e^{B(T)}\leq(1-129/1024)^{-1}=1024/895\), and
\(\kappa\leq v_*+v_*^2=65/4096\). The actual continuum route bound is
\begin{equation}
 r_{route}^2\leq4T\frac{65}{4096}\frac{1024}{895}
                  (C+\epsilon+8c).
 \label{eq:new-numeric-route}
\end{equation}
For an initially coherent route state and one common prepared marker,
the controlled actual-to-reference vector error is at most the largest
branch radius; it is not the sum of two preparation errors.
With the encoder absent and ideal unrecorded routing controls, equal
proper ages give the concrete continuum consequence
\( |\Gamma_{unencoded}-1|\leq2r_{route}\).
The source and shared field thus supply a near-unit-overlap control for
this short closed excursion. This is distinct from the encoded record
experiment below, whose route correlation is paid by a different physical
operation.

Prepare the original marker near age \(T_0/4\), with its retained
preparation-distance budget. After the closed excursion apply the
positive finite route-controlled source pulse
\begin{equation}\begin{gathered}
B_{enc}=\Pi_1\otimes(I_{16}-S)/2,\\
 B_{enc}^2=B_{enc},\\
 \int\lambda(t)dt=\pi,\\
 U_{enc}^{ideal}=\Pi_0\otimes I+\Pi_1\otimes S.
 \label{eq:new-encoder}
\end{gathered}\end{equation}
It acts on all 32 route/source states. It is a quantum controlled
operation: the apparatus does not receive a classical copy of the unknown
route label. Swap anticommutes with the source-lane \(Y\), so the two
reference outputs have opposite \(-Y\) expectations and are decoded by
the same complete source16 POVM. Swap commutes with \(H_f\) and preserves
every source lane. The retained first-energy pulse lemma therefore gives
\[\begin{gathered}
r_{enc}\leq(1+2\pi)\sqrt{2C\tau_{enc}},\\
 \tau_{enc}=2^{-30}.
\end{gathered}\]
The pulse comparison factors out the same lane scalar phase on both route
branches, exactly as in the retained finite-energy lemma. That scalar
remains in the physical action and work account. The physical generator
during the pulse is \(H_0+\lambda B_{enc}\),
with the original even phase retained. Its peak is \(\pi/\tau_{enc}\)
and absolute switch-work payment is at most \(2\pi/\tau_{enc}\).
Those large finite controls are recorded explicitly. This encodes a
route label; it is not an intrinsic gravitational age signal. In
particular the mirror arms have equal proper ages, not the two distinct
ages of the earlier calibration example. Without the encoder their
overlap is given by the shared-field dynamics alone.

The small reference age increment obeys
\(\Omega T\leq\Delta_{up}T<0.01\) in every original lane. Thus the
same declared phase allowance already present in the parent source
decoder covers this reference timing. After ideal encoding the source
field-energy ceiling is unchanged, since \(S\) commutes with \(H_f\).
The actual finite source decoder then obeys its original uniform
form-energy and work bounds on either encoded input. No absolute hard
energy phase needs to be measured by the source-only quadrature.

\subsection{Coherent composition and the noise/work distinction}
At the common return location couple the source to the finite decoder
bit B, the actual nonuniform Chapter IV pointer, retained memory R and
local access bit A. These are the coherent unitary/pulse dilations of
the preceding integrated constructions. They act on marker and apparatus,
with identity on the route register. The complete source, common field,
bit/field/bath correlations and every residual pointer cloud remain.

The scalar response is the actual compact cap background plus the
bit-controlled signed displacement. The latter is a transducer actuator,
not an assignment of an uncomputed intrinsic remote cap mean. A local
station at the common endpoint needs no remote bit-delivery assumption.
For a remote station the physical bit or source signal must arrive through
a specified causal channel before this interaction; no remote record is
being inferred from the local construction.

Use a whole-pipeline hybrid: compare the actual routed and encoded joint
input with the proper-age reference followed by ideal swap encoding;
then apply exactly the same full detector/storage channel to both.
Trace-distance contraction pays \(r_{route}+r_{enc}\) once at the final
readout. It is not used to transfer an unbounded energy moment.
The parent detector moment is recalculated on its reference input using
the unchanged field-energy ceiling and actual decoder work. The actual
encoder energy has its separate, possibly much larger work ceiling.
Both payments remain in the physical energy ledger. A small state error
never licenses assigning the actual encoded state the reference energy.
There is a separate actual-energy envelope. At the closed return,
\(\langle L\rangle\leq E_*e^{B(T)}\), and bounded-pulse energy balance
through encoding and decoding gives
\begin{equation}
 F_{actual,in}\leq2\left[E_*e^{B(T)}
       +W_{enc}^{abs}+W_{dec}^{abs}+c\right].
 \label{eq:new-actual-energy}
\end{equation}
The physical field-gate resource ledger uses this larger ceiling, not
the reference noise moment used in the final probability comparison.
For its actual energy-admissible cap test \(f_t\), put
\(\mu_t=\|f_t/\sqrt\omega\|^2\) and
\(\dot\mu_t^{test}=\|\dot f_t/\sqrt\omega\|^2\).
The conserved pointer momentum second moment \(p_2\) and annihilation
form estimate bound the controlled cap switching work by
\[
 \begin{aligned} W_{cap}^{abs}&\leq2\sqrt{p_2}\int dt\,
  [|\dot h|\sqrt{\mu_t}+|h|\sqrt{\dot\mu_t^{test}}]
       \sqrt{F_{actual}(t)}\\ &\quad+p_2\operatorname{TV}(s).\end{aligned}
\]
This follows from \(|\langle P\phi(f)\rangle|
\leq2\sqrt{p_2\mu_f F_{actual}}\) on the joint state, with no
bit/field independence or fourth moment. For the stabilized gate's
common-form graph, let
\(a(t)G\leq q_t+b(t)\) and
\(|\dot q_t+\dot b(t)|\leq\ell(t)G\). Its actual energy satisfies
\(\mathcal E(t)\leq\mathcal E(0)\exp\int\ell/a\), using
Eq.~\eqref{eq:new-actual-energy} at entry. The retained smooth compact
test/gate and positive stabilizer give finite graph coefficients and
cap-test norms. This is a finite functional ceiling for actual gate work
even if the actual energy greatly exceeds its reference value. Global
background vacuum-ramp and material actuator work remain separately
specified external payments, as in the parent; they are not set to zero
or inferred from the noise ceiling.

Consequently the worst-label final ready error obeys
\begin{equation}
\begin{gathered}
e_{new}^{r}\leq e_{record}^{r}+r_{route}+r_{enc},\\
 e_{all}\leq f+(1-f)e_{new}^{r}.
\end{gathered}
\label{eq:new-complete-error}
\end{equation}
Here the parent record bound already includes preparation, source decode,
pointer loading, finite writing, initialization, hold, aperture and local
read error. Those terms are not charged a second time. The native
operand used here is the parent's conservative \texttt{full\_error}
return; its smaller refined aperture/read display is not silently used
to improve this budget. An encoder area-calibration error \(a_{cal}\)
adds at most \(|a_{cal}|\) by bounded-generator Duhamel. The numerical
example specifies nominal encoder area exactly, just as its nominal
trajectory; it assigns no achieved control calibration. The native
six-lane evaluation certifies \(e_{new}^{r}<0.0016\) with the displayed
chosen controls; \(f\) remains the supplied per-label failure ceiling.
The particular \(f=10^{-4}\) line is an explicit budget, not an achieved
failure measurement.

For the same coherent ensemble, route visibility and this readable
route record obey the integrated tight theorem
\[
 V\leq4\sqrt{pq\,e_{all}(1-e_{all})}.
\]
The marker-only decoder, writing and access operations leave the
unsorted route marginal exactly unchanged. They make an existing route
correlation readable; they do not additionally suppress that marginal
merely because a later observer reads it. Conditional erasure or
route feed-forward is a separately specified experiment.

\subsection{All outcomes, resources and latency}
Retain read results 0 and 1, every preparation/aperture/control failure
flag, and all unobserved recipients in one trace-preserving instrument.
A ready-branch conditional calculation never discards its failed input
weight from the all-attempt ledger. The source decoder and record pulses
remain active finite-duration operations with their original positive
completions and pointer kinetic energy. The new experiment adds the
route work ceiling, encoder capital/peak/switch work and durations
\(T+\tau_{enc}\) to those parent resources. The long analog and logical
holds remain; the source clock certificate applies at source entry, not
through the entire long receiving latency. Microscopic actuator, absolute
length/time and mass calibration remain physical model assignments.

